\clearpage\section*{Solutions EV4-T2}\label{EV4-T2-CORR}
\textbf{MCQ:} 1 D, 2 A, 3 B, 4 A, 5 B, 6 C, 7 A, 8 C, 9 C, 10 D, 11 D, 12 A. Each correct answer: 0.5 point; total 6.
\paragraph{Checks of the fluency questions.} Order of operations; sum of fractions; subtraction of exponents; multiplying coefficient 1.2; division by 0.8; product 1.25 times 0.8; solving an affine equation; vertical increment divided by horizontal; mean; angle divided by 360; complement to 1; substitution in brackets.
\subsection*{Exercise 1 -- A reversed tree}
\begin{enumerate}
\item $P(A\cap B)=\dfrac{1}{5}$; $P(B)=\dfrac{1}{5}+\dfrac{3}{10}=\dfrac{1}{2}$. \textit{Marking: 1 point. Result and basic justification expected; half a point possible for a correct but unfinished approach.}
\item $P_B(A)=\dfrac{2}{5}$; $P(A\cap\overline B)=\dfrac{2}{5}$ so $P_{\overline B}(A)=\dfrac{4}{5}$. \textit{Marking: 2 points. Method or justification: 1; calculation and conclusion: 1.}
\item First level $\dfrac{1}{2}$, $\dfrac{1}{2}$; after B: $\dfrac{2}{5}$ and $\dfrac{3}{5}$; after the complement of B: $\dfrac{4}{5}$ and $\dfrac{1}{5}$. \textit{Marking: 2 points. Method or justification: 1; calculation and conclusion: 1.}
\end{enumerate}
\subsection*{Exercise 2 -- Two optimisation methods}
\begin{enumerate}
\item $f=16-(x-4)^2$; maximum 16 at 4. \textit{Marking: 1 point. Result and basic justification expected; half a point possible for a correct but unfinished approach.}
\item $f\prime=8-2x$, positive before 4 and negative after. \textit{Marking: 2 points. Method or justification: 1; calculation and conclusion: 1.}
\item Endpoints: f(1)=7, f(6)=12; minimum 7. Inequality $(x-4)^2\leq1$: $[3;5]$. \textit{Marking: 2 points. Method or justification: 1; calculation and conclusion: 1.}
\end{enumerate}
\subsection*{Exercise 3 -- Arguments and symmetry}
\begin{enumerate}
\item $A(-\dfrac{\sqrt 3}{2};-\dfrac{1}{2})$, $B(\dfrac{\sqrt 3}{2};-\dfrac{1}{2})$. \textit{Marking: 1 point. Result and basic justification expected; half a point possible for a correct but unfinished approach.}
\item $-\dfrac{3}{4}+\dfrac{1}{4}=-\dfrac{1}{2}$. \textit{Marking: 1 point. Result and basic justification expected; half a point possible for a correct but unfinished approach.}
\item Angle $\dfrac{2\pi}{3}$; $AB=\sqrt3$. \textit{Marking: 2 points. Method or justification: 1; calculation and conclusion: 1.}
\end{enumerate}
\paragraph{Calibration.} Target duration: 20--25 min for the fluency questions, 80--85 min for the exercises, 10--15 min for proofreading. These durations are editorial estimates to be measured in class. Accept any correct method; do not penalise twice a numerical error that is propagated consistently. The proposed marking scheme is not a decision of the team to harmonise marking.
